\section{Experiments}
\label{sec:experiments}

\subsection{Hardware and software}

All development and linear-probe evaluation were performed on an Apple M1
laptop with 8\,GB of unified memory using PyTorch's MPS backend. Pretraining
runs were executed on a single NVIDIA RTX 5090 (32\,GB VRAM, Blackwell
architecture sm\_120) provisioned on the AutoDL cloud platform, with
\texttt{bfloat16} mixed precision and a batch size of 64. The total
pretraining$+$probe pipeline---10{,}000 SSL steps plus the LOSO probe
sweep---runs end-to-end in under 15 minutes on a single consumer GPU. We
provide the complete reproducibility bundle (code, configurations, the six
trained checkpoints used in this paper, and the per-fold CSV outputs) at
\url{https://github.com/lifeinpassion/eeg-lejepa}.

\subsection{Pretraining defaults}

Unless otherwise stated, pretraining uses AdamW
($\beta_1{=}0.9$, $\beta_2{=}0.95$, weight decay $0.05$), initial learning
rate $10^{-3}$ with cosine decay to $10^{-4}$ across training, a linear
warm-up over 30 steps, and gradient clipping at L2-norm 1.0. The SIGReg
weight is $\lambda{=}1.0$ and the number of Cram\'er--Wold projection
slices is $M{=}1024$, both calibrated on a 3-subject 200-step pilot. The
default pretraining corpus is the EEGMMIDB MI runs of subjects 51--100 at
batch size 64 for 10{,}000 steps. For the larger 7\,M-parameter variant
used in the capacity test (\cref{sec:results:capacity}) we lower the
learning rate to $5{\times}10^{-4}$ and extend warm-up to 300 steps, as
required to stabilize training at the higher capacity.

\subsection{Evaluation protocols}

\paragraph{Downstream tasks.}
We probe two EEGMMIDB binary tasks (rest-vs-activity and left-vs-right
motor imagery, runs 4/8/12) and one BCI-IV-2a four-class task (left
hand vs.\ right hand vs.\ both feet vs.\ tongue motor imagery, training
session). For each task and each pretrained checkpoint, we extract
features under all three sources (\texttt{encoder\_mean},
\texttt{predictor\_mean}, \texttt{both\_mean}) and report all three.

\paragraph{Cross-subject LOSO probe.}
For EEGMMIDB evaluation we use leave-one-subject-out cross-validation
across 20 held-out subjects (subjects 1--20); for BCI-IV-2a we use LOSO
across the 9 subjects. The classifier is
\texttt{sklearn.LogisticRegression} ($C{=}1.0$, 2{,}000-iteration cap). We
report per-fold accuracy and per-fold macro one-vs-rest AUC, summarised as
mean and per-fold standard deviation across folds. For the four-class BCI
task we additionally report balanced accuracy and macro-$F_1$.

\paragraph{Supervised baseline.}
For \cref{sec:results:supervised} we train a supervised-from-scratch
controlled comparison: the encoder architecture is identical to the
pretrained one, but with a single linear classification head replacing
the predictor (total $\sim$1.08\,M parameters, smaller than the SSL
pipeline because the predictor is dropped). Training uses cross-entropy
over the same LOSO folds, 500 steps per fold, batch 16, AdamW at
$5{\times}10^{-4}$ with a 50-step linear warm-up, gradient clipping at
1.0. Defensive NaN handling and a try/except around the per-fold AUC
computation are required because the random-initialized classifier
occasionally produces degenerate logits at the first few steps. All other
hyperparameters match the pretrained setup.

\paragraph{No-pretraining control.}
For every probe number reported, we run an identical-architecture model
with random initialization through the same feature-extraction and LOSO
classification pipeline. This control isolates the contribution of SSL
pretraining from any contribution of the architecture or probe protocol.
